\section{Backbone: Sparse U-Net}
\label{sec:backbone}

\subsection{Architecture Overview}

The backbone is a three-level sparse U-Net implemented with
MinkowskiEngine~\cite{choy20194d}.
Figure~\ref{fig:architecture} shows the full architecture including the
FiLM conditioning head.
The total parameter count is approximately 2.5 million.

\begin{figure}[t]
  \centering
  \resizebox{0.98\textwidth}{!}{%
  \begin{tikzpicture}[
    >=Stealth,
    every node/.style={font=\small},
    box/.style={draw, rounded corners=3pt, inner sep=5pt,
                align=center, minimum width=2.6cm, minimum height=0.65cm},
    enc/.style={box, fill=blue!10, draw=blue!40},
    dec/.style={box, fill=teal!10, draw=teal!40},
    film/.style={box, fill=orange!15, draw=orange!55},
    rate/.style={box, fill=orange!10, draw=orange!40},
    sec/.style={font=\small\bfseries},
  ]
  %% ---- Region boxes ----
  \draw[rounded corners=5pt, fill=blue!5, draw=blue!25, dashed]
    (-2.0,-3.6) rectangle (2.0, 0.3);
  \draw[rounded corners=5pt, fill=teal!5, draw=teal!25, dashed]
    (3.0,-3.6) rectangle (7.0, 0.3);
  \draw[rounded corners=5pt, fill=orange!5, draw=orange!30, dashed]
    (7.5, 0.6) rectangle (19.0, 2.1);
  %% ---- Region labels ----
  \node[sec, blue!65,  rotate=90]  at (-2.45, -1.65) {Encoder};
  \node[sec, teal!70,  rotate=-90] at (7.45,  -1.65) {Decoder};
  \node[sec, orange!70] at (13.2, 1.97) {Rate Conditioning};
  %% ---- Nodes ----
  \node[box, fill=gray!10] (P)   at (0,    0.85) {Input $\mathbf{P}$\\{\scriptsize 3-ch $xyz$}};
  \node[enc] (e1) at (0,   -0.8) {E$_1$: ResBlock, 32ch};
  \node[enc] (e2) at (0,   -2.0) {E$_2$: ResBlock, 64ch};
  \node[enc] (e3) at (0,   -3.0) {E$_3$: ResBlock, 128ch};
  \node[dec] (d3) at (5.0, -3.0) {D$_3$: ResBlock, 128ch};
  \node[dec] (d2) at (5.0, -2.0) {D$_2$: ResBlock, 64ch};
  \node[dec] (d1) at (5.0, -0.8) {D$_1$: ResBlock, 32ch};
  \node[film](fh) at (5.0,  1.3)
    {FiLM Head\\{\scriptsize $\mathbf{h}' = \boldsymbol{\gamma}(r)\odot\mathbf{h}+\boldsymbol{\beta}(r)$}};
  \node[box, fill=gray!10] (out) at (5.0, 2.55)
    {$\hat{\mathbf{P}} = \mathbf{P} + \Delta$};
  \node[rate](gb)  at (9.5,  1.3) {$\boldsymbol{\gamma},\boldsymbol{\beta}\in\mathbb{R}^{32}$};
  \node[rate](mlp) at (13.2, 1.3) {MLP\ {\scriptsize $1{\to}64{\to}64$}};
  \node[rate](bpp) at (17.0, 1.3) {$r{=}\log_2(\mathrm{bpp})$};
  %% ---- Info text (fills white space right of enc/dec, below rate cond) ----
  \node[anchor=north west, font=\small] at (7.8, 0.1) {%
    \begin{tabular}{@{}l@{\enspace}l@{}}
      \textbf{Rate MLP:}    & $[\boldsymbol{\gamma};\boldsymbol{\beta}] =
                               W_2\,\sigma(W_1 r + b_1) + b_2,\quad \sigma=\mathrm{ReLU}$ \\[4pt]
      \textbf{Modulation:}  & $\mathbf{h}' = \boldsymbol{\gamma}(r)\odot\mathbf{h}+\boldsymbol{\beta}(r),
                               \quad \boldsymbol{\gamma},\boldsymbol{\beta}\in\mathbb{R}^{32}$ \\[4pt]
      \textbf{Output head:} & SparseConv $1^3$;\quad $\Delta = \tanh(\cdot)\times 5$ \\[4pt]
      \textbf{Result:}      & $\hat{\mathbf{P}} = \mathbf{P} + \Delta$\quad (point count preserved)
    \end{tabular}%
  };
  %% ---- Arrows ----
  \draw[->] (P)  -- (e1);
  \draw[->] (e1) -- node[right]{\scriptsize $\times2\!\downarrow$} (e2);
  \draw[->] (e2) -- node[right]{\scriptsize $\times2\!\downarrow$} (e3);
  \draw[->] (e3) -- (d3);
  \draw[->] (d3) -- node[right]{\scriptsize $\div2\!\uparrow$} (d2);
  \draw[->] (d2) -- node[right]{\scriptsize $\div2\!\uparrow$} (d1);
  \draw[->] (d1) -- node[right]{\scriptsize 32-ch} (fh);
  \draw[->] (fh) -- (out);
  \draw[->, dashed, blue!55, thick]
    (e1.east) -- node[above]{\scriptsize feat.\ concat} (d1.west);
  \draw[->, dashed, blue!55, thick]
    (e2.east) -- node[above]{\scriptsize feat.\ concat} (d2.west);
  \draw[->] (bpp) -- (mlp);
  \draw[->] (mlp) -- node[above]{\scriptsize split 32+32} (gb);
  \draw[->] (gb.west) -- (fh.east);
  \end{tikzpicture}%
  }
  \caption{%
    Architecture of the rate-conditioned post-processing network.
    The rate-invariant sparse U-Net backbone (Encoder + Decoder) processes
    the decoded point cloud.
    Skip connections (dashed) concatenate encoder features to decoder inputs
    at matching resolutions.
    The FiLM head applies per-channel affine modulation conditioned on
    $r = \log_2(\mathrm{bpp})$; the output head produces bounded displacement vectors
    $\Delta$ added to the input to yield the refined cloud.
  }
  \label{fig:architecture}
\end{figure}

\subsection{Encoder Path}

The encoder applies three levels of processing:
\begin{align*}
  E_1 &: \text{3-ch} \to 32\text{ ch},\quad \text{SubConv}_{3^3} \text{ ResBlock},\quad \text{stride }1 \\
  E_2 &: 32 \to 64\text{ ch},\quad \text{StrideConv}_{2^3} + \text{ResBlock},\quad \text{stride }2 \\
  E_3 &: 64 \to 128\text{ ch},\quad \text{StrideConv}_{2^3} + \text{ResBlock},\quad \text{stride }2
\end{align*}

At each encoder level $\ell$, a ResBlock with two $3^3$ submanifold sparse
convolutions (Section~\ref{ssec:sparse-conv}) processes the features at
the current resolution.
At $E_2$ and $E_3$, a stride-2 sparse convolution precedes the ResBlock,
halving the spatial resolution and doubling the channel count.

After three levels, the receptive field at the finest resolution is
$1 + 2 \times (3-1) + 4 \times (3-1) = 13$ voxels per side, providing
a $13^3 = 2197$ voxel neighbourhood for each active coordinate at the
bottleneck level.
This is sufficient to capture local surface geometry (edges, ridges) at the
scale of fine features in the 8iVFB clouds.

\subsection{Decoder Path and Skip Connections}

The decoder mirrors the encoder with transposed (upsampling) stride-2
convolutions and skip connections from the encoder:
\begin{align*}
  D_3 &: (128 + 128) \to 128\text{ ch},\quad \text{ResBlock},\quad \text{stride }1 \\
  D_2 &: \text{TrConv}(\times2\uparrow) + (64 + 64) \to 64\text{ ch},\quad \text{ResBlock} \\
  D_1 &: \text{TrConv}(\times2\uparrow) + (32 + 32) \to 32\text{ ch},\quad \text{ResBlock}
\end{align*}

At each decoder level, the upsampled features from the level below are
concatenated with the skip connection from the matching encoder level.
The concatenated features have twice the channel count of either branch alone,
and a ResBlock reduces them back to the target channel count.

\paragraph{Why skip connections matter here.}
The displacement prediction is a per-point, high-resolution task.
Without skip connections, the only geometric information available to the
decoder at fine resolution is what propagates up through the bottleneck,
at 4$\times$ downsampled resolution.
Skip connections from $E_1$ and $E_2$ provide access to the fine-resolution
coordinate features at the original spatial scale, which are critical for
predicting displacement direction and magnitude per voxel.

\subsection{Rate Invariance of the Backbone}

The backbone weights are shared across all compression rates.
This is a deliberate design choice:
\begin{itemize}
  \item The geometric correction task shares common structure across rates:
    the network must identify surface discontinuities and displace decoded
    points toward the original surface regardless of the rate.
  \item Sharing weights provides cross-rate regularisation: the backbone
    is exposed to the same correction task at different distortion scales,
    encouraging it to learn a rate-invariant geometric representation.
  \item Rate-specific adaptation is handled entirely by the FiLM head
    (Section~\ref{sec:film-head}), which modulates the backbone output at
    inference time.
\end{itemize}

% ============================================================
